\section{RQ1: Overall Performance}

\subsection{主表：Instruments 上当前最重要的 run 地图}

\begin{table}[H]
\centering
\footnotesize
\setlength{\tabcolsep}{4pt}
\begin{tabular}{p{1.9cm} p{3.8cm} p{2.3cm} c c c}
\toprule
Group & Variant & Best ckpt & Epoch & NDCG@10 & HR@50 \\
\midrule
Baseline & SFT495 & checkpoint-495 & -- & 0.0823 & 0.1844 \\
Baseline & rule-only rerun & checkpoint-2997 & 1.802 & 0.0960 & 0.1681 \\
Dynamic & dynamic gather-fix & checkpoint-2997 & 1.802 & 0.0936 & 0.1855 \\
Dynamic & dynamic sid-only & checkpoint-2394 & 1.805 & 0.0921 & 0.1830 \\
Dynamic & ranking dynamic & checkpoint-1665 & 1.001 & 0.0911 & 0.1822 \\
Fixed & fixed taskfix & checkpoint-2997 & 1.802 & 0.0931 & 0.1941 \\
Fixed & fixed taskfix sid-only & checkpoint-2652 & 2.000 & 0.0945 & 0.1935 \\
Historical & old fixed mixed-single & checkpoint-3326 & 2.000 & 0.0953 & 0.1938 \\
Ablation & dynamic max1 & checkpoint-1332 & 0.801 & 0.0934 & 0.1905 \\
CE & hintce-2 & checkpoint-2664 & 1.602 & 0.0931 & 0.1951 \\
CE & hintce-3 & checkpoint-1665 & 1.001 & 0.0945 & 0.1985 \\
Hint ablation & single-hint mixed & checkpoint-2664 & 1.602 & 0.0948 & 0.1958 \\
Dual-task & fixed dual-task & checkpoint-3012 & 2.000 & 0.0939 & 0.1905 \\
\bottomrule
\end{tabular}
\caption{RQ1 的主结果地图：Instruments 上当前已经同步回来的关键 run。}
\end{table}

\begin{figure}[H]
\centering
\includegraphics[width=0.96\textwidth]{rl-seven-way-main-curves.png}
\caption{主线七线曲线图，用来定位 RQ1 的整体版图。}
\end{figure}

\begin{figure}[H]
\centering
\includegraphics[width=0.78\textwidth]{rl-best-ndcg10-vs-hr50-scatter.png}
\caption{best \texttt{NDCG@10} 对应的 top-10 / coverage trade-off。}
\end{figure}

\subsection{No-hint baseline 与 hint-enabled frontier}

\begin{table}[H]
\centering
\small
\setlength{\tabcolsep}{4pt}
\begin{tabular}{p{4.0cm} p{2.4cm} c c c c}
\toprule
Variant & Best ckpt & NDCG@10 & HR@10 & NDCG@50 & HR@50 \\
\midrule
SFT495 & checkpoint-495 & 0.0823 & 0.1094 & 0.0985 & 0.1844 \\
rule-only rerun & checkpoint-2997 & 0.0960 & 0.1179 & 0.1070 & 0.1681 \\
dynamic gather-fix & checkpoint-2997 & 0.0936 & 0.1169 & 0.1083 & 0.1855 \\
fixed taskfix sid-only & checkpoint-2652 & 0.0945 & 0.1205 & 0.1103 & 0.1935 \\
\bottomrule
\end{tabular}
\caption{RQ1 最核心的四条线：SFT、plain rule-only、canonical adaptive hinting、clean fixed first-token hint。}
\end{table}

这张表给出正文后续所有 RQ 的共同起点：

\begin{itemize}
\item \texttt{rule-only} 仍然是当前最强的无 hint top-10 baseline，但 coverage 损失很明显。
\item \texttt{dynamic gather-fix} 说明一旦在训练里显式引入 prefix hint，模型就能把 \texttt{HR@50} 拉回到 SFT 以上，而不是只会做 headline overfit。
\item \texttt{fixed taskfix sid-only} 进一步把 trade-off 推到更右上，说明 prefix hint 的收益并不只停留在 adaptive family。
\end{itemize}

\subsection{RQ1 的简化读法}

如果只想记住这一节的结论，可以压缩成三点：

\begin{enumerate}
\item plain \texttt{rule-only} 仍然定义了“无 hint 时 top-10 能推到多高”的 baseline。
\item canonical \texttt{dynamic gather-fix} 说明 hint-conditioned training signal 是有效的，因为它没有像 \texttt{rule-only} 那样用 coverage 去换 headline。
\item clean fixed family、尤其是 \texttt{fixed taskfix sid-only} 与后面的 CE / partial-hint 候选，构成了当前真正需要继续比较的 performance frontier。
\end{enumerate}
